\input format.tex

\usepackage{graphicx}
\graphicspath{{cores/}}

\usepackage{environ}
\usepackage{colortbl,array,booktabs}
\usepackage{tabularx}

\colorlet{TablaBordeSuperior}{topcolor}
\colorlet{TablaBordeInferior}{topcolor}
\colorlet{TablaCentroSuperior}{blue!1}
\colorlet{TablaCentroInferior}{blue!20}
\colorlet{FuenteCabeceraTabla}{white}

\newcolumntype{M}[1]{>{\centering\arraybackslash}m{#1}}

\tcbset{rtab/.style={
freelance,
frame code={
 \path[top color=topcolor,bottom color=topcolor]
   ([yshift=-#1*(\baselineskip+2pt)]interior.north west) --
   ([yshift=-#1*(\baselineskip+2pt)]interior.north east) {[sharp corners]--
    ([yshift=3pt]interior.north east) --
    ([yshift=3pt]interior.north west)} -- cycle;

  },
interior code={},
 }
}

\newcommand\fuentecabecera[1]{\textcolor{black}{\textbf{#1}}}

\NewEnviron{RCtable}[1]
  {%
    \begin{tcolorbox}[#1]\BODY\end{tcolorbox}%
  }


\begin{document}

\vspace*{4mm}
%% 各章节
\setlength{\arrayrulewidth}{1pt}
\fontsize{9.3pt}{11pt}\selectfont
\color{gray2}
\noindent
{\fontsize{9pt}{11pt}\selectfont\bf 致病菌检测结果}
\begin{tctabularx}{tabularx={m{4.8cm}<{\centering}m{5.2cm}<{\centering}m{4.55cm}<{\centering}}}
&&
\\[-6pt]
\cellcolor{topcolor} \raisebox{2.614pt}{\color{gray0}\bf 检测项目} &
\cellcolor{topcolor} \raisebox{2.614pt}{\color{gray0}\bf 健康状态提示} &
\cellcolor{topcolor} \raisebox{2.614pt}{\color{gray0}\bf 检测结果评价}
\end{tctabularx}

\vspace*{-4.25mm}
\fontsize{8pt}{11pt}\selectfont
\arrayrulecolor{gray2}
\begin{longtable}{m{4.8cm}m{5.2cm}<{\centering}m{0cm}@{}m{4.61cm}<{\centering}}
\hline
\parbox[c]{\hsize}{\vskip7pt {\lantxh 缓慢爱格士氏菌\\可能导致菌血症、椎间盘炎、脑膜炎等} \vskip7pt} & \parbox[c]{\hsize}{\vskip7pt\centerline{\raisebox{-1.5ex}{\includegraphics[width=5cm,height=0.55cm]{bar.pdf}}}\vskip7pt}  &
\hspace*{-1.51cm}\raisebox{-0.45ex}{\includegraphics{cry.pdf}}
 & \begin{minipage}{4.60cm}\begin{center}{{\color{red}\lantxh 高{\\ \bahao 可能增加菌血症、椎间盘炎、脑膜炎、脑脓肿等的风险}} }\end{center} \end{minipage} \\
\hline
\parbox[c]{\hsize}{\vskip7pt {\lantxh 胶胨罗斯氏菌\\可能导致败血症、脑膜炎、肺炎、菌血症等} \vskip7pt} & \parbox[c]{\hsize}{\vskip7pt\centerline{\raisebox{-1.5ex}{\includegraphics[width=5cm,height=0.55cm]{bar.pdf}}}\vskip7pt}  &
\hspace*{-1.51cm}\raisebox{-0.45ex}{\includegraphics{cry.pdf}}
 & \begin{minipage}{4.60cm}\begin{center}{{\color{red}\lantxh 高{\\ \bahao 可能增加败血症、脑膜炎、肺炎、菌血症等的风险}} }\end{center} \end{minipage} \\
\hline
\parbox[c]{\hsize}{\vskip7pt {\lantxh 副流感嗜血杆菌\\可导致脑膜炎、咽炎、心内膜炎、肺炎、败血症等} \vskip7pt} & \parbox[c]{\hsize}{\vskip7pt\centerline{\raisebox{-1.5ex}{\includegraphics[width=5cm,height=0.55cm]{bar.pdf}}}\vskip7pt}  &
\hspace*{-1.51cm}\raisebox{-0.45ex}{\includegraphics{cry.pdf}}
 & \begin{minipage}{4.60cm}\begin{center}{{\color{red}\lantxh 高{\\ \bahao 可能增加脑膜炎、咽炎、心内膜炎、肺炎、败血症等的风险}} }\end{center} \end{minipage} \\
\hline
\parbox[c]{\hsize}{\vskip7pt {\lantxh 咽峡炎链球菌\\属于人体正常菌群一部分，但也可能造成肝脓肿、脑脓肿、菌血症等} \vskip7pt} & \parbox[c]{\hsize}{\vskip7pt\centerline{\raisebox{-1.5ex}{\includegraphics[width=5cm,height=0.55cm]{bar.pdf}}}\vskip7pt}  &
\hspace*{-1.51cm}\raisebox{-0.45ex}{\includegraphics{cry.pdf}}
 & \begin{minipage}{4.60cm}\begin{center}{{\color{red}\lantxh 高{\\ \bahao 可能增加肝脓肿、脑脓肿、菌血症等的风险}} }\end{center} \end{minipage} \\
\hline
\multicolumn{3}{l}{}\\
\multicolumn{3}{l}{\includegraphics[scale=1]{greencircle.pdf} {\fontsize{9.3pt}{11pt}\selectfont \bf 绿色表示健康}   \includegraphics[scale=1]{orangecircle.pdf}  {\fontsize{9.3pt}{11pt}\selectfont \bf 橙色表示需要关注}  \includegraphics[scale=1]{redcircle.pdf} {\fontsize{9.3pt}{11pt}\selectfont \bf 红色表示有风险}}
\end{longtable}

%\vspace*{8mm}
\vspace*{6mm}
\begin{spacing}{1.25}
\fontsize{9.3pt}{11pt}\selectfont
\begin{LRaside}[.8]{结果分析}
\noindent
\includegraphics[width=\linewidth]{result_fenbu.pdf}
\asidebreak %
1、您肠道内的有害菌含量偏高，会产生毒素等有害物质，可能会引起口臭、感染、腹泻、肠炎、便秘等。有害菌数量过高还可能影响您的心情和食欲，扰乱内分泌，降低机体免疫力，增加疾病风险。但同时您肠道内的有益菌含量也较高，数量较多的有益菌能够起到拮抗作用，降低有害菌可能对您造成的危害。有害菌在您免疫力低下时，可能导致相关疾病，建议您规律作息，注意饮食健康，预防潜在的健康风险。\\
2、您的肠道内
检测出4种致病菌，其中缓慢爱格士氏菌、胶胨罗斯氏菌、副流感嗜血杆菌、咽峡炎链球菌含量高于参考人群，当身体免疫力较强时可能不会出现不适感。但请注意，当您的免疫力较弱时，初期可能引起腹泻、腹痛、感染等病症，长期可能导致上表中所述疾病。建议您持续监测，必要时请到正规医院咨询专科医生，遵从医生的建议做进一步检查。
\end{LRaside}
\end{spacing}

\end{document}
